\title{2022 年高考数学 全国乙卷-理科}
\date{}
\maketitle

\section*{一、选择题}
本题共12小题，每小题5分，共60分．在每小题给出的四个选项中，只有一项是符合题目要求的．

\begin{question}[points=5]
设全集$U = \{1, 2, 3, 4, 5\}$，集合$M$满足$\complement_U M = \{1, 3\}$，则（\qquad）
\begin{choices}
  \item $2 \in M$
  \item $3 \in M$
  \item $4 \notin M$
  \item $5 \notin M$
\end{choices}
\begin{solution}
\textbf{答案：}A

\textbf{分析：}先写出集合$M$,然后逐项验证即可

\textbf{详解：}由题知$M = \{2, 4, 5\}$,对比选项知，A正确, BCD错误
\end{solution}
\end{question}

\begin{question}[points=5]
已知$z = 1 - 2\mathrm{i}$，且$z + a\bar{z} + b = 0$，其中$a$，$b$为实数，则（\qquad）
\begin{choices}
  \item $a = 1, b = -2$
  \item $a = -1, b = 2$
  \item $a = 1, b = 2$
  \item $a = -1, b = -2$
\end{choices}
\begin{solution}
\textbf{答案：}A

\textbf{分析：}先算出$\bar{z}$,再代入计算,实部与虚部都为零解方程组即可

\textbf{详解：}$\bar{z} = 1 + 2\mathrm{i}$\\$z + a\bar{z} + b = 1 - 2\mathrm{i} + a(1 + 2\mathrm{i}) + b = (1 + a + b) + (2a - 2)\mathrm{i}$\\由$z + a\bar{z} + b = 0$,得$\begin{cases} 1 + a + b = 0 \\ 2a - 2 = 0 \end{cases}$,即$\begin{cases} a = 1 \\ b = -2 \end{cases}$
\end{solution}
\end{question}

\begin{question}[points=5]
已知向量$\vec{a}, \vec{b}$满足$|\vec{a}| = 1, |\vec{b}| = \sqrt{3}, |\vec{a} - 2\vec{b}| = 3$，则$\vec{a} \cdot \vec{b} =$（\qquad）
\begin{choices}
  \item $-2$
  \item $-1$
  \item $1$
  \item $2$
\end{choices}
\begin{solution}
\textbf{答案：}C

\textbf{分析：}根据给定模长，利用向量的数量积运算求解即可.

\textbf{详解：}解：∵$|\vec{a} - 2\vec{b}|^2 = |\vec{a}|^2 - 4\vec{a} \cdot \vec{b} + 4|\vec{b}|^2$，\\又∵$|\vec{a}| = 1, |\vec{b}| = \sqrt{3}, |\vec{a} - 2\vec{b}| = 3$，\\∴$9 = 1 - 4\vec{a} \cdot \vec{b} + 4 \times 3 = 13 - 4\vec{a} \cdot \vec{b}$，\\∴$\vec{a} \cdot \vec{b} = 1$
\end{solution}
\end{question}

\begin{question}[points=5]
嫦娥二号卫星在完成探月任务后，继续进行深空探测，成为我国第一颗环绕太阳飞行的人造行星，为研究嫦娥二号绕日周期与地球绕日周期的比值，用到数列$\{b_n\}$：$b_1 = 1 + \frac{1}{\alpha_1}$，$b_2 = 1 + \frac{1}{\alpha_1 + \frac{1}{\alpha_2}}$，$b_3 = 1 + \frac{1}{\alpha_1 + \frac{1}{\alpha_2 + \frac{1}{\alpha_3}}}$，…，依此类推，其中$\alpha_k \in \mathbb{N}^* (k = 1, 2, \cdots)$．则（\qquad）
\begin{choices}
  \item $b_1 < b_5$
  \item $b_3 < b_8$
  \item $b_6 < b_2$
  \item $b_4 < b_7$
\end{choices}
\begin{solution}
\textbf{答案：}D

\textbf{分析：}根据$\alpha_k \in \mathbb{N}^* (k = 1, 2, \dots)$，再利用数列$\{b_n\}$与$\alpha_k$的关系判断$\{b_n\}$中各项的大小，即可求解.

\textbf{详解：}解：因为$\alpha_k \in \mathbb{N}^* (k = 1, 2, \cdots)$，\\所以$\alpha_1 < \alpha_1 + \frac{1}{\alpha_2}$，$\alpha_1 + \frac{1}{\alpha_2} > \frac{1}{\alpha_1 + \frac{1}{\alpha_2}}$，得到$b_1 > b_2$，\\同理$\alpha_1 + \frac{1}{\alpha_2} > \alpha_1 + \frac{1}{\alpha_2 + \frac{1}{\alpha_3}}$，可得$b_2 < b_3$，$b_1 > b_3$\\又因为$\frac{1}{\alpha_2} > \frac{1}{\alpha_2 + \frac{1}{\alpha_3 + \frac{1}{\alpha_4}}}$，$\alpha_1 + \frac{1}{\alpha_2 + \frac{1}{\alpha_3}} < \alpha_1 + \frac{1}{\alpha_2 + \frac{1}{\alpha_3 + \frac{1}{\alpha_4}}}$，故$b_2 < b_4$，$b_3 > b_4$；\\以此类推，可得$b_1 > b_3 > b_5 > b_7 > \dots$，$b_7 > b_8$，故A错误；\\$b_1 > b_7 > b_8$，故B错误；\\$\frac{1}{\alpha_2} > \frac{1}{\alpha_2 + \frac{1}{\alpha_3 + \dots \frac{1}{\alpha_6}}}$，得$b_2 < b_6$，故C错误；\\$\alpha_1 + \frac{1}{\alpha_2 + \frac{1}{\alpha_3 + \frac{1}{\alpha_4}}} > \alpha_1 + \frac{1}{\alpha_2 + \dots \frac{1}{\alpha_6 + \frac{1}{\alpha_7}}}$，得$b_4 < b_7$，故D正确.
\end{solution}
\end{question}

\begin{question}[points=5]
设$F$为抛物线$C: y^2 = 4x$的焦点，点$A$在$C$上，点$B(3, 0)$，若$|AF| = |BF|$，则$|AB| =$（\qquad）
\begin{choices}
  \item $2$
  \item $2\sqrt{2}$
  \item $3$
  \item $3\sqrt{2}$
\end{choices}
\begin{solution}
\textbf{答案：}B

\textbf{分析：}根据抛物线上的点到焦点和准线的距离相等，从而求得点$A$的横坐标，进而求得点$A$坐标，即可得到答案.

\textbf{详解：}由题意得，$F(1,0)$，则$|AF| = |BF| = 2$，\\即点$A$到准线$x=-1$的距离为2，所以点$A$的横坐标为$-1+2=1$，\\不妨设点$A$在$x$轴上方，代入得，$A(1,2)$，\\所以$|AB| = \sqrt{(3-1)^2 + (0-2)^2} = 2\sqrt{2}$.
\end{solution}
\end{question}

\begin{question}[points=5]
执行下边的程序框图，输出的$n =$（\qquad）
\begin{choices}
  \item $3$
  \item $4$
  \item $5$
  \item $6$
\end{choices}
\begin{solution}
\textbf{答案：}B

\textbf{分析：}根据框图循环计算即可.

\textbf{详解：}执行第一次循环，$b = b + 2a = 1 + 2 = 3$，\\$a = b - a = 3 - 1 = 2, n = n + 1 = 2$，\\$\frac{b^2}{a^2} - 2 = \frac{9}{4} - 2 = \frac{1}{4} > 0.01$；\\执行第二次循环，$b = b + 2a = 3 + 4 = 7$，\\$a = b - a = 7 - 2 = 5, n = n + 1 = 3$，\\$\frac{b^2}{a^2} - 2 = \frac{49}{25} - 2 = \frac{1}{25} > 0.01$；\\执行第三次循环，$b = b + 2a = 7 + 10 = 17$，\\$a = b - a = 17 - 5 = 12, n = n + 1 = 4$，\\$\frac{b^2}{a^2} - 2 = \frac{289}{144} - 2 = \frac{1}{144} < 0.01$，此时输出$n=4$.
\end{solution}
\end{question}

\begin{question}[points=5]
在正方体$ABCD-A_1B_1C_1D_1$中，$E$，$F$分别为$AB$，$BC$的中点，则（\qquad）
\begin{choices}
  \item 平面$B_1EF \perp$平面$BDD_1$
  \item 平面$B_1EF \perp$平面$A_1BD$
  \item 平面$B_1EF \parallel$平面$A_1AC$
  \item 平面$B_1EF \parallel$平面$A_1C_1D$
\end{choices}
\begin{solution}
\textbf{答案：}A

\textbf{分析：}证明$EF \perp$平面$BDD_1$，即可判断A；如图，以点$D$为原点，建立空间直角坐标系，设$AB=2$，分别求出平面$B_1EF$，$A_1BD$，$A_1C_1D$的法向量，根据法向量的位置关系，即可判断BCD.

\textbf{详解：}解：在正方体$ABCD-A_1B_1C_1D_1$中，\\$AC \perp BD$且$DD_1 \perp$平面$ABCD$，\\又$EF \subset$平面$ABCD$，所以$EF \perp DD_1$，\\因为$E,F$分别为$AB,BC$的中点，\\所以$EF \parallel AC$，所以$EF \perp BD$，\\又$BD \cap DD_1 = D$，\\所以$EF \perp$平面$BDD_1$，\\又$EF \subset$平面$B_1EF$，\\所以平面$B_1EF \perp$平面$BDD_1$，故A正确；\\如图，以点$D$为原点，建立空间直角坐标系，设$AB=2$，\\则$B_1(2,2,2), E(2,1,0), F(1,2,0), B(2,2,0), A_1(2,0,2), A(2,0,0), C(0,2,0), C_1(0,2,2)$，\\$\overrightarrow{EF}=(-1,1,0), \overrightarrow{EB_1}=(0,1,2), \overrightarrow{DB}=(2,2,0), \overrightarrow{DA_1}=(2,0,2), \overrightarrow{AA_1}=(0,0,2), \overrightarrow{AC}=(-2,2,0), \overrightarrow{A_1C_1}=(-2,2,0)$，\\设平面$B_1EF$的法向量为$\vec{m}=(x_1,y_1,z_1)$，\\则有$\begin{cases} \vec{m}\cdot\overrightarrow{EF}=-x_1+y_1=0 \\ \vec{m}\cdot\overrightarrow{EB_1}=y_1+2z_1=0 \end{cases}$，可取$\vec{m}=(2,2,-1)$，\\同理可得平面$A_1BD$的法向量为$\vec{n_1}=(1,-1,-1)$，\\平面$A_1AC$的法向量为$\vec{n_2}=(1,1,0)$，\\平面$A_1C_1D$的法向量为$\vec{n_3}=(1,1,-1)$，\\则$\vec{m}\cdot\vec{n_1}=2-2+1=1\neq0$，\\所以平面$B_1EF$与平面$A_1BD$不垂直，故B错误；\\因为$\vec{m}$与$\vec{n_2}$不平行，\\所以平面$B_1EF$与平面$A_1AC$不平行，故C错误；\\因为$\vec{m}$与$\vec{n_3}$不平行，\\所以平面$B_1EF$与平面$A_1C_1D$不平行，故D错误.
\end{solution}
\end{question}

\begin{question}[points=5]
已知等比数列$\{a_n\}$的前3项和为168，$a_2 - a_5 = 42$，则$a_6 =$（\qquad）
\begin{choices}
  \item $14$
  \item $12$
  \item $6$
  \item $3$
\end{choices}
\begin{solution}
\textbf{答案：}D

\textbf{分析：}设等比数列$\{a_n\}$的公比为$q, q\neq0$，易得$q\neq1$，根据题意求出首项与公比，再根据等比数列的通项即可得解.

\textbf{详解：}解：设等比数列$\{a_n\}$的公比为$q, q\neq0$，\\若$q=1$，则$a_2-a_5=0$，与题意矛盾，\\所以$q\neq1$，\\则$\begin{cases} a_1+a_2+a_3 = \frac{a_1(1-q^3)}{1-q} = 168 \\ a_2 - a_5 = a_1q - a_1q^4 = 42 \end{cases}$，解得$\begin{cases} a_1=96 \\ q=\frac{1}{2} \end{cases}$，\\所以$a_6 = a_1q^5 = 3$.
\end{solution}
\end{question}

\begin{question}[points=5]
已知球$O$的半径为1，四棱锥的顶点为$O$，底面的四个顶点均在球$O$的球面上，则当该四棱锥的体积最大时，其高为（\qquad）
\begin{choices}
  \item $\frac{1}{3}$
  \item $\frac{1}{2}$
  \item $\frac{\sqrt{3}}{3}$
  \item $\frac{\sqrt{2}}{2}$
\end{choices}
\begin{solution}
\textbf{答案：}C

\textbf{分析：}先证明当四棱锥的顶点$O$到底面$ABCD$所在小圆距离一定时，底面$ABCD$面积最大值为$2r^2$，进而得到四棱锥体积表达式，再利用均值定理去求四棱锥体积的最大值，从而得到当该四棱锥的体积最大时其高的值.

\textbf{详解：}设该四棱锥底面为四边形$ABCD$，四边形$ABCD$所在小圆半径为$r$，设四边形$ABCD$对角线夹角为$\alpha$，\\则$S_{ABCD} = \frac{1}{2}\cdot AC\cdot BD\cdot\sin\alpha \leq \frac{1}{2}\cdot AC\cdot BD \leq \frac{1}{2}\cdot 2r\cdot 2r = 2r^2$（当且仅当四边形$ABCD$为正方形时等号成立）\\即当四棱锥的顶点$O$到底面$ABCD$所在小圆距离一定时，底面$ABCD$面积最大值为$2r^2$\\又$r^2 + h^2 = 1$\\则$V_{O-ABCD} = \frac{1}{3}\cdot 2r^2\cdot h = \frac{2}{3}r\cdot r\cdot 2h^2 \leq \frac{2}{3}\left(\frac{r^2 + r^2 + 2h^2}{3}\right)^3 = \frac{4\sqrt{3}}{27}$\\当且仅当$r^2 = 2h^2$即$h = \frac{\sqrt{3}}{3}$时等号成立.
\end{solution}
\end{question}

\begin{question}[points=5]
某棋手与甲、乙、丙三位棋手各比赛一盘，各盘比赛结果相互独立．已知该棋手与甲、乙、丙比赛获胜的概率分别为$p_1, p_2, p_3$，且$p_3 > p_2 > p_1 > 0$．记该棋手连胜两盘的概率为$p$，则（\qquad）
\begin{choices}
  \item $p$与该棋手和甲、乙、丙的比赛次序无关
  \item 该棋手在第二盘与甲比赛，$p$最大
  \item 该棋手在第二盘与乙比赛，$p$最大
  \item 该棋手在第二盘与丙比赛，$p$最大
\end{choices}
\begin{solution}
\textbf{答案：}D

\textbf{分析：}该棋手连胜两盘，则第二盘为必胜盘.分别求得该棋手在第二盘与甲比赛且连胜两盘的概率$p_{\text{甲}}$；该棋手在第二盘与乙比赛且连胜两盘的概率$p_{\text{乙}}$；该棋手在第二盘与丙比赛且连胜两盘的概率$p_{\text{丙}}$.并对三者进行比较即可解决

\textbf{详解：}该棋手连胜两盘，则第二盘为必胜盘，\\记该棋手在第二盘与甲比赛，且连胜两盘的概率为$p_{\text{甲}}$\\则$p_{\text{甲}} = 2(1-p_2)p_1p_3 + 2p_2p_1(1-p_3) = 2p_1(p_2+p_3) - 4p_1p_2p_3$\\记该棋手在第二盘与乙比赛，且连胜两盘的概率为$p_{\text{乙}}$\\则$p_{\text{乙}} = 2(1-p_1)p_2p_3 + 2p_1p_2(1-p_3) = 2p_2(p_1+p_3) - 4p_1p_2p_3$\\记该棋手在第二盘与丙比赛，且连胜两盘的概率为$p_{\text{丙}}$\\则$p_{\text{丙}} = 2(1-p_1)p_3p_2 + 2p_1p_3(1-p_2) = 2p_3(p_1+p_2) - 4p_1p_2p_3$\\则$p_{\text{甲}} - p_{\text{乙}} = 2p_1(p_2+p_3) - 4p_1p_2p_3 - [2p_2(p_1+p_3) - 4p_1p_2p_3] = 2(p_1-p_2)p_3 < 0$\\$p_{\text{乙}} - p_{\text{丙}} = 2p_2(p_1+p_3) - 4p_1p_2p_3 - [2p_3(p_1+p_2) - 4p_1p_2p_3] = 2(p_2-p_3)p_1 < 0$\\即$p_{\text{甲}} < p_{\text{乙}}$，$p_{\text{乙}} < p_{\text{丙}}$，\\则该棋手在第二盘与丙比赛，$p$最大.
\end{solution}
\end{question}

\begin{question}[points=5]
双曲线$C$的两个焦点为$F_1, F_2$，以$C$的实轴为直径的圆记为$D$，过$F_1$作$D$的切线与$C$的两支交于$M, N$两点，且$\cos\angle F_1NF_2 = \frac{3}{5}$，则$C$的离心率为（\qquad）
\begin{choices}
  \item $\frac{\sqrt{5}}{2}$
  \item $\frac{3}{2}$
  \item $\frac{\sqrt{13}}{2}$
  \item $\frac{\sqrt{17}}{2}$
\end{choices}
\begin{solution}
\textbf{答案：}C

\textbf{分析：}依题意不妨设双曲线焦点在$x$轴，设过$F_1$作圆$D$的切线切点为$G$，可判断$N$在双曲线的右支，设$\angle F_1NF_2 = \alpha$，$\angle F_2F_1N = \beta$，即可求出$\sin\alpha$，$\sin\beta$，$\cos\beta$，在$\triangle F_2F_1N$中由$\sin\angle F_1F_2N = \sin(\alpha+\beta)$求出$\sin\angle F_1F_2N$，再由正弦定理求出$NF_1$，$NF_2$，最后根据双曲线的定义得到$2b=3a$，即可得解；

\textbf{详解：}解：依题意不妨设双曲线焦点在$x$轴，设过$F_1$作圆$D$的切线切点为$G$，\\所以$OG \perp NF_1$，因为$\cos\angle F_1NF_2 = \frac{3}{5} > 0$，所以$N$在双曲线的右支，\\所以$OG = a$，$OF_1 = c$，$GF_1 = b$，设$\angle F_1NF_2 = \alpha$，$\angle F_2F_1N = \beta$，\\由$\cos\angle F_1NF_2 = \frac{3}{5}$，即$\cos\alpha = \frac{3}{5}$，则$\sin\alpha = \frac{4}{5}$，$\sin\beta = \frac{a}{c}$，$\cos\beta = \frac{b}{c}$，\\在$\triangle F_2F_1N$中，$\sin\angle F_1F_2N = \sin(\pi-\alpha-\beta) = \sin(\alpha+\beta)$\\$= \sin\alpha\cos\beta + \cos\alpha\sin\beta = \frac{4}{5}\cdot\frac{b}{c} + \frac{3}{5}\cdot\frac{a}{c} = \frac{3a+4b}{5c}$，\\由正弦定理得$\frac{2c}{\sin\alpha} = \frac{NF_2}{\sin\beta} = \frac{NF_1}{\sin\angle F_1F_2N} = \frac{5c}{2}$，\\所以$NF_1 = \frac{5c}{2}\sin\angle F_1F_2N = \frac{5c}{2}\cdot\frac{3a+4b}{5c} = \frac{3a+4b}{2}$，$NF_2 = \frac{5c}{2}\sin\beta = \frac{5c}{2}\cdot\frac{a}{c} = \frac{5a}{2}$\\又$NF_1 - NF_2 = \frac{3a+4b}{2} - \frac{5a}{2} = \frac{4b-2a}{2} = 2a$，\\所以$2b = 3a$，即$\frac{b}{a} = \frac{3}{2}$，\\所以双曲线的离心率$e = \frac{c}{a} = \sqrt{1+\frac{b^2}{a^2}} = \frac{\sqrt{13}}{2}$.
\end{solution}
\end{question}

\begin{question}[points=5]
已知函数$f(x), g(x)$的定义域均为$\mathbb{R}$，且$f(x) + g(2-x) = 5$，$g(x) - f(x-4) = 7$．若$y = g(x)$的图像关于直线$x=2$对称，$g(2) = 4$，则$\sum\limits_{k=1}^{22} f(k) =$（\qquad）
\begin{choices}
  \item $-21$
  \item $-22$
  \item $-23$
  \item $-24$
\end{choices}
\begin{solution}
\textbf{答案：}D

\textbf{分析：}根据对称性和已知条件得到$f(x) + f(x-2) = -2$，从而得到$f(3)+f(5)+\cdots+f(21) = -10$，$f(4)+f(6)+\cdots+f(22) = -10$，然后根据条件得到$f(2)$的值，再由题意得到$g(3)=6$从而得到$f(1)$的值即可求解.

\textbf{详解：}因为$y=g(x)$的图像关于直线$x=2$对称，\\所以$g(2-x) = g(x+2)$，\\因为$g(x) - f(x-4) = 7$，所以$g(x+2) - f(x-2) = 7$，即$g(x+2) = 7 + f(x-2)$，\\因为$f(x) + g(2-x) = 5$，所以$f(x) + g(x+2) = 5$，\\代入得$f(x) + [7 + f(x-2)] = 5$，即$f(x) + f(x-2) = -2$，\\所以$f(3)+f(5)+\cdots+f(21) = (-2)\times5 = -10$，\\$f(4)+f(6)+\cdots+f(22) = (-2)\times5 = -10$，\\因为$f(x) + g(2-x) = 5$，所以$f(0) + g(2) = 5$，即$f(0) = 1$，所以$f(2) = -2 - f(0) = -3$，\\因为$g(x) - f(x-4) = 7$，所以$g(x+4) - f(x) = 7$，又因为$f(x) + g(2-x) = 5$，\\联立得，$g(2-x) + g(x+4) = 12$，\\所以$y=g(x)$的图像关于点$(3,6)$中心对称，因为函数$g(x)$的定义域为$\mathbb{R}$，\\所以$g(3)=6$\\因为$f(x) + g(x+2) = 5$，所以$f(1) = 5 - g(3) = -1$，\\所以$\sum\limits_{k=1}^{22} f(k) = f(1)+f(2)+[f(3)+f(5)+\cdots+f(21)]+[f(4)+f(6)+\cdots+f(22)] = -1 -3 -10 -10 = -24$.
\end{solution}
\end{question}

\section*{二、填空题}
本题共4小题，每小题5分，共20分．

\begin{question}[points=5]
从甲、乙等5名同学中随机选3名参加社区服务工作，则甲、乙都入选的概率为．
\fillin[{$\frac{3}{10}$}]
\begin{solution}
\textbf{答案：}$\frac{3}{10}$

\textbf{分析：}根据古典概型计算即可

\textbf{详解：}从5名同学中随机选3名的方法数为$\mathrm{C}_5^3 = 10$\\甲、乙都入选的方法数为$\mathrm{C}_3^1 = 3$，所以甲、乙都入选的概率$P = \frac{3}{10}$
\end{solution}
\end{question}

\begin{question}[points=5]
过四点$(0,0),(4,0),(-1,1),(4,2)$中的三点的一个圆的方程为．
\fillin[{$(x-2)^2 + (y-3)^2 = 13$或$(x-2)^2 + (y-1)^2 = 5$或$\left(x-\frac{4}{3}\right)^2 + \left(y-\frac{7}{3}\right)^2 = \frac{65}{9}$或$\left(x-\frac{8}{5}\right)^2 + (y-1)^2 = \frac{169}{25}$}]
\begin{solution}
\textbf{答案：}$(x-2)^2 + (y-3)^2 = 13$或$(x-2)^2 + (y-1)^2 = 5$或$\left(x-\frac{4}{3}\right)^2 + \left(y-\frac{7}{3}\right)^2 = \frac{65}{9}$或$\left(x-\frac{8}{5}\right)^2 + (y-1)^2 = \frac{169}{25}$

\textbf{分析：}设圆的方程为$x^2+y^2+Dx+Ey+F=0$，根据所选点的坐标，得到方程组，解得即可；

\textbf{详解：}解：依题意设圆的方程为$x^2+y^2+Dx+Ey+F=0$，\\若过$(0,0),(4,0),(-1,1)$，则$\begin{cases} F=0 \\ 16+4D+F=0 \\ 1+1-D+E+F=0 \end{cases}$，解得$\begin{cases} D=-4 \\ E=-6 \\ F=0 \end{cases}$，\\所以圆的方程为$x^2+y^2-4x-6y=0$，即$(x-2)^2+(y-3)^2=13$；\\若过$(0,0),(4,0),(4,2)$，则$\begin{cases} F=0 \\ 16+4D+F=0 \\ 16+4+4D+2E+F=0 \end{cases}$，解得$\begin{cases} D=-4 \\ E=-2 \\ F=0 \end{cases}$，\\所以圆的方程为$x^2+y^2-4x-2y=0$，即$(x-2)^2+(y-1)^2=5$；\\若过$(0,0),(4,2),(-1,1)$，则$\begin{cases} F=0 \\ 1+1-D+E+F=0 \\ 16+4+4D+2E+F=0 \end{cases}$，解得$\begin{cases} D=-\frac{8}{3} \\ E=-\frac{14}{3} \\ F=0 \end{cases}$，\\所以圆的方程为$x^2+y^2-\frac{8}{3}x-\frac{14}{3}y=0$，即$\left(x-\frac{4}{3}\right)^2+\left(y-\frac{7}{3}\right)^2=\frac{65}{9}$；\\若过$(-1,1),(4,0),(4,2)$，则$\begin{cases} 1+1-D+E+F=0 \\ 16+4D+F=0 \\ 16+4+4D+2E+F=0 \end{cases}$，解得$\begin{cases} D=-\frac{16}{5} \\ E=-2 \\ F=-\frac{16}{5} \end{cases}$，\\所以圆的方程为$x^2+y^2-\frac{16}{5}x-2y-\frac{16}{5}=0$，即$\left(x-\frac{8}{5}\right)^2+(y-1)^2=\frac{169}{25}$；
\end{solution}
\end{question}

\begin{question}[points=5]
记函数$f(x) = \cos(\omega x + \varphi)$ $(\omega > 0, 0 < \varphi < \pi)$的最小正周期为$T$，若$f(T) = \frac{\sqrt{3}}{2}$，且$x = \frac{\pi}{9}$为$f(x)$的零点，则$\omega$的最小值为．
\fillin[{$3$}]
\begin{solution}
\textbf{答案：}$3$

\textbf{分析：}首先表示出$T$，根据$f(T)=\frac{\sqrt{3}}{2}$求出$\varphi$，再根据$x=\frac{\pi}{9}$为函数的零点，即可求出$\omega$的取值，从而得解；

\textbf{详解：}解：因为$f(x)=\cos(\omega x+\varphi)$，$(\omega>0,0<\varphi<\pi)$\\所以最小正周期$T=\frac{2\pi}{\omega}$，因为$f(T)=\cos\left(\omega\cdot\frac{2\pi}{\omega}+\varphi\right)=\cos(2\pi+\varphi)=\cos\varphi=\frac{\sqrt{3}}{2}$，\\又$0<\varphi<\pi$，所以$\varphi=\frac{\pi}{6}$，即$f(x)=\cos\left(\omega x+\frac{\pi}{6}\right)$，\\又$x=\frac{\pi}{9}$为$f(x)$的零点，所以$\frac{\pi}{9}\omega+\frac{\pi}{6}=\frac{\pi}{2}+k\pi, k\in\mathbb{Z}$，解得$\omega=3+9k, k\in\mathbb{Z}$，\\因为$\omega>0$，所以当$k=0$时$\omega_{\min}=3$；
\end{solution}
\end{question}

\begin{question}[points=5]
已知$x=x_1$和$x=x_2$分别是函数$f(x)=2a^x - ex^2$（$a>0$且$a\neq1$）的极小值点和极大值点．若$x_1<x_2$，则$a$的取值范围是．
\fillin[{$\left(\frac{1}{e},1\right)$}]
\begin{solution}
\textbf{答案：}$\left(\frac{1}{e},1\right)$

\textbf{分析：}由$x_1,x_2$分别是函数$f(x)=2a^x-ex^2$的极小值点和极大值点，可得$x\in(-\infty,x_1)\cup(x_2,+\infty)$时，$f'(x)<0$，$x\in(x_1,x_2)$时，$f'(x)>0$，再分$a>1$和$0<a<1$两种情况讨论，方程$2\ln a\cdot a^x - 2ex = 0$的两个根为$x_1,x_2$，即函数$y=\ln a\cdot a^x$与函数$y=ex$的图象有两个不同的交点，构造函数$g(x)=\ln a\cdot a^x$，根据导数的结合意义结合图象即可得出答案.

\textbf{详解：}解：$f'(x)=2\ln a\cdot a^x - 2ex$，\\因为$x_1,x_2$分别是函数$f(x)=2a^x-ex^2$的极小值点和极大值点，\\所以函数$f(x)$在$(-\infty,x_1)$和$(x_2,+\infty)$上递减，在$(x_1,x_2)$上递增，\\所以当$x\in(-\infty,x_1)\cup(x_2,+\infty)$时，$f'(x)<0$，当$x\in(x_1,x_2)$时，$f'(x)>0$，\\若$a>1$时，\\当$x<0$时，$2\ln a\cdot a^x>0, 2ex<0$，\\则此时$f'(x)>0$，与前面矛盾，\\故$a>1$不符合题意，\\若$0<a<1$时，\\则方程$2\ln a\cdot a^x - 2ex = 0$的两个根为$x_1,x_2$，\\即方程$\ln a\cdot a^x = ex$的两个根为$x_1,x_2$，\\即函数$y=\ln a\cdot a^x$与函数$y=ex$的图象有两个不同的交点，\\令$g(x)=\ln a\cdot a^x$，则$g'(x)=\ln^2 a\cdot a^x, 0<a<1$，\\设过原点且与函数$y=g(x)$的图象相切的直线的切点为$(x_0,\ln a\cdot a^{x_0})$，\\则切线的斜率为$g'(x_0)=\ln^2 a\cdot a^{x_0}$，\\故切线方程为$y-\ln a\cdot a^{x_0}=\ln^2 a\cdot a^{x_0}(x-x_0)$，\\则有$-\ln a\cdot a^{x_0} = -x_0\ln^2 a\cdot a^{x_0}$，\\解得$x_0=\frac{1}{\ln a}$，\\则切线的斜率为$\ln^2 a\cdot a^{\frac{1}{\ln a}} = e\ln^2 a$，\\因为函数$y=\ln a\cdot a^x$与函数$y=ex$的图象有两个不同的交点，\\所以$e\ln^2 a < e$，解得$\frac{1}{e}<a<e$，\\又$0<a<1$，所以$\frac{1}{e}<a<1$，\\综上所述，$a$的范围为$\left(\frac{1}{e},1\right)$.
\end{solution}
\end{question}

\section*{三、解答题}
解答应写出文字说明、证明过程或演算步骤．第17~21题为必考题，每个试题考生都必须作答．第22、23题为选考题，考生根据要求作答．

\begin{question}[points=12]
记$\triangle ABC$的内角$A,B,C$的对边分别为$a,b,c$，已知$\sin C\sin(A-B) = \sin B\sin(C-A)$．

(1)证明：$2a^2 = b^2 + c^2$；

(2)若$a=5,\cos A = \frac{25}{31}$，求$\triangle ABC$的周长．
\begin{solution}
\textbf{答案：}(1)见解析 (2)14

\textbf{分析：}(1)利用两角差的正弦公式化简，再根据正弦定理和余弦定理化角为边，从而即可得证；(2)根据(1)的结论结合余弦定理求出$bc$，从而可求得$b+c$，即可得解.

\textbf{详解：}【小问1详解】\\证明：因为$\sin C\sin(A-B)=\sin B\sin(C-A)$，\\所以$\sin C\sin A\cos B - \sin C\sin B\cos A = \sin B\sin C\cos A - \sin B\sin A\cos C$，\\所以$ac\cdot\frac{a^2+c^2-b^2}{2ac} - 2bc\cdot\frac{b^2+c^2-a^2}{2bc} = -ab\cdot\frac{a^2+b^2-c^2}{2ab}$，\\即$\frac{a^2+c^2-b^2}{2} - (b^2+c^2-a^2) = -\frac{a^2+b^2-c^2}{2}$，\\所以$2a^2 = b^2 + c^2$；\\【小问2详解】\\解：因为$a=5,\cos A=\frac{25}{31}$，\\由(1)得$b^2+c^2=50$，\\由余弦定理可得$a^2 = b^2+c^2 - 2bc\cos A$，\\则$50 - \frac{50}{31}bc = 25$，\\所以$bc = \frac{31}{2}$，\\故$(b+c)^2 = b^2+c^2+2bc = 50+31=81$，\\所以$b+c=9$，\\所以$\triangle ABC$的周长为$a+b+c=14$.
\end{solution}
\end{question}

\begin{question}[points=12]
如图，四面体$ABCD$中，$AD\perp CD, AD=CD, \angle ADB = \angle BDC$，$E$为$AC$的中点．

(1)证明：平面$BED \perp$平面$ACD$；

(2)设$AB=BD=2, \angle ACB=60^{\circ}$，点$F$在$BD$上，当$\triangle AFC$的面积最小时，求$CF$与平面$ABD$所成的角的正弦值．

\begin{center}
  \includegraphics[width=.38\linewidth]{question-18-1.jpg}
\end{center}
\begin{solution}
\textbf{答案：}(1)证明过程见解析 (2)$\frac{4\sqrt{3}}{7}$

\textbf{分析：}(1)根据已知关系证明$\triangle ABD \cong \triangle CBD$，得到$AB=CB$，结合等腰三角形三线合一得到垂直关系，结合面面垂直的判定定理即可证明；(2)根据勾股定理逆用得到$BE\perp DE$，从而建立空间直角坐标系，结合线面角的运算法则进行计算即可.

\textbf{详解：}【小问1详解】\\因为$AD=CD$，$E$为$AC$的中点，所以$AC\perp DE$；\\在$\triangle ABD$和$\triangle CBD$中，因为$AD=CD, \angle ADB=\angle CDB, DB=DB$，\\所以$\triangle ABD \cong \triangle CBD$，所以$AB=CB$，又因为$E$为$AC$的中点，所以$AC\perp BE$；\\又因为$DE, BE\subset$平面$BED$，$DE\cap BE=E$，所以$AC\perp$平面$BED$，\\因为$AC\subset$平面$ACD$，所以平面$BED\perp$平面$ACD$.\\【小问2详解】\\连接$EF$，由(1)知，$AC\perp$平面$BED$，因为$EF\subset$平面$BED$，\\所以$AC\perp EF$，所以$S_{\triangle AFC}=\frac{1}{2}AC\cdot EF$，\\当$EF\perp BD$时，$EF$最小，即$\triangle AFC$的面积最小.\\因为$\triangle ABD\cong\triangle CBD$，所以$CB=AB=2$，\\又因为$\angle ACB=60^{\circ}$，所以$\triangle ABC$是等边三角形，\\因为$E$为$AC$的中点，所以$AE=EC=1$，$BE=\sqrt{3}$，\\因为$AD\perp CD$，所以$DE=\frac{1}{2}AC=1$，\\在$\triangle DEB$中，$DE^2+BE^2=BD^2$，所以$BE\perp DE$.\\以$E$为坐标原点建立如图所示的空间直角坐标系$E-xyz$，\\则$A(1,0,0), B(0,\sqrt{3},0), D(0,0,1)$，所以$\overrightarrow{AD}=(-1,0,1), \overrightarrow{AB}=(-1,\sqrt{3},0)$，\\设平面$ABD$的一个法向量为$\vec{n}=(x,y,z)$，\\则$\begin{cases} \vec{n}\cdot\overrightarrow{AD} = -x+z=0 \\ \vec{n}\cdot\overrightarrow{AB} = -x+\sqrt{3}y=0 \end{cases}$，取$y=\sqrt{3}$，则$\vec{n}=(3,\sqrt{3},3)$，\\又因为$C(-1,0,0), F\left(0,\frac{\sqrt{3}}{4},\frac{3}{4}\right)$，所以$\overrightarrow{CF}=\left(1,\frac{\sqrt{3}}{4},\frac{3}{4}\right)$，\\所以$\cos\langle\vec{n},\overrightarrow{CF}\rangle = \frac{\vec{n}\cdot\overrightarrow{CF}}{|\vec{n}||\overrightarrow{CF}|} = \frac{6}{\sqrt{21}\times\frac{7}{4}} = \frac{4\sqrt{3}}{7}$，\\设$CF$与平面$ABD$所成的角的正弦值为$\theta\left(0\leq\theta\leq\frac{\pi}{2}\right)$，\\所以$\sin\theta = |\cos\langle\vec{n},\overrightarrow{CF}\rangle| = \frac{4\sqrt{3}}{7}$，\\所以$CF$与平面$ABD$所成的角的正弦值为$\frac{4\sqrt{3}}{7}$.
\end{solution}
\end{question}

\begin{question}[points=12]
某地经过多年的环境治理，已将荒山改造成了绿水青山．为估计一林区某种树木的总材积量，随机选取了10棵这种树木，测量每棵树的根部横截面积（单位：$m^2$）和材积量（单位：$m^3$），得到如下数据：



| 样本号$i$ | 1 | 2 | 3 | 4 | 5 | 6 | 7 | 8 | 9 | 10 | 总和 |

|-----------|---|---|---|---|---|---|---|---|---|---|------|

| 根部横截面积$x_i$ | 0.04 | 0.06 | 0.04 | 0.08 | 0.08 | 0.05 | 0.05 | 0.07 | 0.07 | 0.06 | 0.6 |

| 材积量$y_i$ | 0.25 | 0.40 | 0.22 | 0.54 | 0.51 | 0.34 | 0.36 | 0.46 | 0.42 | 0.40 | 3.9 |



并计算得$\sum\limits_{i=1}^{10} x_i^2 = 0.038, \sum\limits_{i=1}^{10} y_i^2 = 1.6158, \sum\limits_{i=1}^{10} x_i y_i = 0.2474$．

(1)估计该林区这种树木平均一棵的根部横截面积与平均一棵的材积量；

(2)求该林区这种树木的根部横截面积与材积量的样本相关系数（精确到0.01）；

(3)现测量了该林区所有这种树木的根部横截面积，并得到所有这种树木的根部横截面积总和为$186m^2$．已知树木的材积量与其根部横截面积近似成正比．利用以上数据给出该林区这种树木的总材积量的估计值．

附：相关系数$r=\frac{\sum\limits_{i=1}^n (x_i-\bar{x})(y_i-\bar{y})}{\sqrt{\sum\limits_{i=1}^n (x_i-\bar{x})^2 \sum\limits_{i=1}^n (y_i-\bar{y})^2}}$，$\sqrt{1.896} \approx 1.377$．
\begin{solution}
\textbf{答案：}(1) $0.06m^2$；$0.39m^3$
(2) $0.97$
(3) $1209m^3$

\textbf{分析：}(1)计算出样本的一棵根部横截面积的平均值及一棵材积量平均值，即可估计该林区这种树木平均一棵的根部横截面积与平均一棵的材积量；(2)代入题给相关系数公式去计算即可求得样本的相关系数值；(3)依据树木的材积量与其根部横截面积近似成正比，列方程即可求得该林区这种树木的总材积量的估计值．

\textbf{详解：}【小问1详解】\\样本中10棵这种树木的根部横截面积的平均值$\bar{x} = \frac{0.6}{10} = 0.06$\\样本中10棵这种树木的材积量的平均值$\bar{y} = \frac{3.9}{10} = 0.39$\\据此可估计该林区这种树木平均一棵的根部横截面积为$0.06m^2$，平均一棵的材积量为$0.39m^3$\\【小问2详解】\\$r = \frac{\sum\limits_{i=1}^{10} (x_i-\bar{x})(y_i-\bar{y})}{\sqrt{\sum\limits_{i=1}^{10} (x_i-\bar{x})^2 \sum\limits_{i=1}^{10} (y_i-\bar{y})^2}} = \frac{\sum\limits_{i=1}^{10} x_i y_i - 10\bar{x}\bar{y}}{\sqrt{(\sum\limits_{i=1}^{10} x_i^2 - 10\bar{x}^2)(\sum\limits_{i=1}^{10} y_i^2 - 10\bar{y}^2)}}$\\$= \frac{0.2474 - 10\times0.06\times0.39}{\sqrt{(0.038 - 10\times0.06^2)(1.6158 - 10\times0.39^2)}} = \frac{0.0134}{\sqrt{0.0001896}} \approx \frac{0.0134}{0.01377} \approx 0.97$\\则$r \approx 0.97$\\【小问3详解】\\设该林区这种树木的总材积量的估计值为$Y\,m^3$，\\又已知树木的材积量与其根部横截面积近似成正比，\\可得$\frac{0.06}{0.39} = \frac{186}{Y}$，解之得$Y=1209m^3$．\\则该林区这种树木的总材积量估计为$1209m^3$
\end{solution}
\end{question}

\begin{question}[points=12]
已知椭圆$E$的中心为坐标原点，对称轴为$x$轴、$y$轴，且过$A(0,-2), B\left(\frac{3}{2},-1\right)$两点．

(1)求$E$的方程；

(2)设过点$P(1,-2)$的直线交$E$于$M,N$两点，过$M$且平行于$x$轴的直线与线段$AB$交于点$T$，点$H$满足$\overrightarrow{MT}=\overrightarrow{TH}$．证明：直线$HN$过定点．
\begin{solution}
\textbf{答案：}(1) $\frac{y^2}{4}+\frac{x^2}{3}=1$
(2) $(0,-2)$

\textbf{分析：}(1)将给定点代入设出的方程求解即可；(2)设出直线方程，与椭圆C的方程联立，分情况讨论斜率是否存在，即可得解.

\textbf{详解：}【小问1详解】\\解：设椭圆E的方程为$mx^2+ny^2=1$，过$A(0,-2), B\left(\frac{3}{2},-1\right)$，\\则$\begin{cases} 4n=1 \\ \frac{9}{4}m+n=1 \end{cases}$，解得$m=\frac{1}{3}, n=\frac{1}{4}$，\\所以椭圆E的方程为：$\frac{y^2}{4}+\frac{x^2}{3}=1$.\\【小问2详解】\\$A(0,-2), B(\frac{3}{2},-1)$，所以$AB: y+2 = \frac{2}{3}x$，\\①若过点$P(1,-2)$的直线斜率不存在，直线$x=1$.代入$\frac{x^2}{3}+\frac{y^2}{4}=1$，\\可得$M\left(1,\frac{2\sqrt{6}}{3}\right), N\left(1,-\frac{2\sqrt{6}}{3}\right)$，代入AB方程$y=\frac{2}{3}x-2$，可得\\$T\left(\sqrt{6}+3,\frac{2\sqrt{6}}{3}\right)$，由$\overrightarrow{MT}=\overrightarrow{TH}$得到$H\left(2\sqrt{6}+5,\frac{2\sqrt{6}}{3}\right)$.求得HN方程：$y=\left(2-\frac{2\sqrt{6}}{3}\right)x-2$，过点$(0,-2)$.\\②若过点$P(1,-2)$的直线斜率存在，设$kx-y-(k+2)=0, M(x_1,y_1), N(x_2,y_2)$.\\联立$\begin{cases} kx-y-(k+2)=0 \\ \frac{x^2}{3}+\frac{y^2}{4}=1 \end{cases}$，得$(3k^2+4)x^2 - 6k(2+k)x + 3k(k+4)=0$，\\可得$\begin{cases} x_1+x_2 = \frac{6k(2+k)}{3k^2+4} \\ x_1x_2 = \frac{3k(4+k)}{3k^2+4} \end{cases}$，$\begin{cases} y_1+y_2 = \frac{-8(2+k)}{3k^2+4} \\ y_1y_2 = \frac{4(4+4k-2k^2)}{3k^2+4} \end{cases}$，\\且$x_1y_2+x_2y_1 = \frac{-24k}{3k^2+4}$ (*)\\联立$\begin{cases} y=y_1 \\ y=\frac{2}{3}x-2 \end{cases}$，可得$T\left(\frac{3y_1}{2}+3, y_1\right), H(3y_1+6-x_1, y_1)$.\\可求得此时$HN: y-y_2 = \frac{y_1-y_2}{3y_1+6-x_1-x_2}(x-x_2)$，\\将$(0,-2)$代入整理得$2(x_1+x_2)-6(y_1+y_2)+x_1y_2+x_2y_1-3y_1y_2-12=0$，\\将(*)代入，得$24k+12k^2+96+48k-24k-48-48k+24k^2-36k^2-48=0$，显然成立，\\综上，可得直线HN过定点$(0,-2)$.
\end{solution}
\end{question}

\begin{question}[points=12]
已知函数$f(x)=\ln(1+x)+axe^{-x}$．

(1)当$a=1$时，求曲线$y=f(x)$在点$(0,f(0))$处的切线方程；

(2)若$f(x)$在区间$(-1,0),(0,+\infty)$各恰有一个零点，求$a$的取值范围．
\begin{solution}
\textbf{答案：}(1) $y=2x$
(2) $(-\infty,-1)$

\textbf{分析：}(1)先算出切点,再求导算出斜率即可；(2)求导,对$a$分类讨论,对$x$分$(-1,0),(0,+\infty)$两部分研究

\textbf{详解：}【小问1详解】\\$f(x)$的定义域为$(-1,+\infty)$\\当$a=1$时,$f(x)=\ln(1+x)+\frac{x}{e^x}, f(0)=0$,所以切点为$(0,0)$\\$f'(x)=\frac{1}{1+x}+\frac{1-x}{e^x}, f'(0)=2$,所以切线斜率为2\\所以曲线$y=f(x)$在点$(0,f(0))$处的切线方程为$y=2x$\\【小问2详解】\\$f(x)=\ln(1+x)+\frac{ax}{e^x}$\\$f'(x)=\frac{1}{1+x}+\frac{a(1-x)}{e^x}=\frac{e^x+a(1-x^2)}{(1+x)e^x}$\\设$g(x)=e^x+a(1-x^2)$\\$1^\circ$若$a>0$,当$x\in(-1,0), g(x)=e^x+a(1-x^2)>0$,即$f'(x)>0$\\所以$f(x)$在$(-1,0)$上单调递增,$f(x)<f(0)=0$\\故$f(x)$在$(-1,0)$上没有零点,不合题意\\$2^\circ$若$-1\leq a\leq 0$,当$x\in(0,+\infty)$,则$g'(x)=e^x-2ax>0$\\所以$g(x)$在$(0,+\infty)$上单调递增所以$g(x)>g(0)=1+a\geq0$,即$f'(x)>0$\\所以$f(x)$在$(0,+\infty)$上单调递增,$f(x)>f(0)=0$\\故$f(x)$在$(0,+\infty)$上没有零点,不合题意\\$3^\circ$若$a<-1$\\(1)当$x\in(0,+\infty)$,则$g'(x)=e^x-2ax>0$,所以$g(x)$在$(0,+\infty)$上单调递增\\$g(0)=1+a<0, g(1)=e>0$\\所以存在$m\in(0,1)$,使得$g(m)=0$,即$f'(m)=0$\\当$x\in(0,m), f'(x)<0, f(x)$单调递减\\当$x\in(m,+\infty), f'(x)>0, f(x)$单调递增\\所以当$x\in(0,m), f(x)<f(0)=0$\\当$x\to+\infty, f(x)\to+\infty$\\所以$f(x)$在$(m,+\infty)$上有唯一零点\\又$(0,m)$没有零点,即$f(x)$在$(0,+\infty)$上有唯一零点\\(2)当$x\in(-1,0), g(x)=e^x+a(1-x^2)$\\设$h(x)=g'(x)=e^x-2ax$\\$h'(x)=e^x-2a>0$\\所以$g'(x)$在$(-1,0)$单调递增\\$g'(-1)=\frac{1}{e}+2a<0, g'(0)=1>0$\\所以存在$n\in(-1,0)$,使得$g'(n)=0$\\当$x\in(-1,n), g'(x)<0, g(x)$单调递减\\当$x\in(n,0), g'(x)>0, g(x)$单调递增$g(x)<g(0)=1+a<0$\\又$g(-1)=\frac{1}{e}>0$\\所以存在$t\in(-1,n)$,使得$g(t)=0$,即$f'(t)=0$\\当$x\in(-1,t), f(x)$单调递增,当$x\in(t,0), f(x)$单调递减\\有$x\to-1, f(x)\to-\infty$\\而$f(0)=0$,所以当$x\in(t,0), f(x)>0$\\所以$f(x)$在$(-1,t)$上有唯一零点,$(t,0)$上无零点\\即$f(x)$在$(-1,0)$上有唯一零点\\所以$a<-1$,符合题意\\所以若$f(x)$在区间$(-1,0),(0,+\infty)$各恰有一个零点,求$a$的取值范围为$(-\infty,-1)$
\end{solution}
\end{question}

\section*{选考题}
共10分．请考生在第22、23题中任选一题作答．如果多做，则按所做的第一题计分．

\begin{question}[points=10]
[选修4-4：坐标系与参数方程]

在直角坐标系$xOy$中，曲线$C$的参数方程为$\begin{cases} x = 3\cos 2t \\ y = 2\sin t \end{cases}$（$t$为参数），以坐标原点为极点，$x$轴正半轴为极轴建立极坐标系，已知直线$l$的极坐标方程为$\rho\sin\left(\theta+\frac{\pi}{3}\right)+m=0$．

(1)写出$l$的直角坐标方程；

(2)若$l$与$C$有公共点，求$m$的取值范围．
\begin{solution}
\textbf{答案：}(1) $\sqrt{3}x+y+2m=0$
(2) $-\frac{19}{12}\leq m\leq \frac{5}{2}$

\textbf{分析：}(1)根据极坐标与直角坐标的互化公式处理即可；(2)联立$l$与$C$的方程，采用换元法处理，根据新设$a$的取值范围求解$m$的范围即可.

\textbf{详解：}【小问1详解】\\因为$l$：$\rho\sin\left(\theta+\frac{\pi}{3}\right)+m=0$，所以$\rho\cdot\frac{1}{2}\sin\theta + \rho\cdot\frac{\sqrt{3}}{2}\cos\theta + m = 0$，\\又因为$\rho\sin\theta=y, \rho\cos\theta=x$，所以化简为$\frac{1}{2}y+\frac{\sqrt{3}}{2}x+m=0$，\\整理得$l$的直角坐标方程：$\sqrt{3}x+y+2m=0$\\【小问2详解】\\联立$l$与$C$的方程，即将$x=3\cos2t, y=2\sin t$代入$\sqrt{3}x+y+2m=0$中，可得$3\sqrt{3}\cos2t+2\sin t+2m=0$，\\所以$3\sqrt{3}(1-2\sin^2t)+2\sin t+2m=0$，\\化简为$-6\sqrt{3}\sin^2t+2\sin t+3\sqrt{3}+2m=0$，\\要使$l$与$C$有公共点，则$2m=6\sqrt{3}\sin^2t-2\sin t-3\sqrt{3}$有解，\\令$\sin t=a$，则$a\in[-1,1]$，令$f(a)=6\sqrt{3}a^2-2a-3\sqrt{3}$，$(-1\leq a\leq 1)$，\\对称轴为$a=\frac{1}{6\sqrt{3}}$，开口向上，\\所以$f(a)_{\max}=f(-1)=6\sqrt{3}+2-3\sqrt{3}=3\sqrt{3}+2$，\\$f(a)_{\min}=f\left(\frac{1}{6\sqrt{3}}\right)=6\sqrt{3}\cdot\frac{1}{108}-\frac{2}{6\sqrt{3}}-3\sqrt{3}=\frac{\sqrt{3}}{18}-\frac{\sqrt{3}}{9}-3\sqrt{3}=-\frac{19\sqrt{3}}{6}$，\\所以$-\frac{19\sqrt{3}}{6}\leq 2m\leq 3\sqrt{3}+2$\\$m$的取值范围为$-\frac{19\sqrt{3}}{12}\leq m\leq \frac{3\sqrt{3}}{2}+1$?? 注意原解析答案？
\end{solution}
\end{question}

\begin{question}[points=10]
[选修4-5：不等式选讲]

已知$a,b,c$都是正数，且$a^{\frac{3}{2}}+b^{\frac{3}{2}}+c^{\frac{3}{2}}=1$，证明：

(1)$abc\leq\frac{1}{9}$；

(2)$\frac{a}{b+c}+\frac{b}{a+c}+\frac{c}{a+b}\leq\frac{1}{2\sqrt{abc}}$．
\begin{solution}
\textbf{答案：}见解析

\textbf{分析：}(1)利用三元均值不等式即可证明；(2)利用基本不等式及不等式的性质证明即可．

\textbf{详解：}【小问1详解】\\证明：因为$a>0,b>0,c>0$，则$a^{\frac{3}{2}}>0,b^{\frac{3}{2}}>0,c^{\frac{3}{2}}>0$，\\所以$a^{\frac{3}{2}}+b^{\frac{3}{2}}+c^{\frac{3}{2}}\geq 3\sqrt[3]{a^{\frac{3}{2}}b^{\frac{3}{2}}c^{\frac{3}{2}}}$，\\即$(abc)^{\frac{1}{2}}\leq\frac{1}{3}$，所以$abc\leq\frac{1}{9}$，当且仅当$a^{\frac{3}{2}}=b^{\frac{3}{2}}=c^{\frac{3}{2}}$即$a=b=c=\sqrt[3]{\frac{1}{9}}$时取等号．\\【小问2详解】\\证明：因为$a>0,b>0,c>0$，\\所以$b+c\geq 2\sqrt{bc}$，$a+c\geq 2\sqrt{ac}$，$a+b\geq 2\sqrt{ab}$，\\所以$\frac{a}{b+c}\leq\frac{a}{2\sqrt{bc}}=\frac{a^{\frac{3}{2}}}{2\sqrt{abc}}$，$\frac{b}{a+c}\leq\frac{b}{2\sqrt{ac}}=\frac{b^{\frac{3}{2}}}{2\sqrt{abc}}$，$\frac{c}{a+b}\leq\frac{c}{2\sqrt{ab}}=\frac{c^{\frac{3}{2}}}{2\sqrt{abc}}$，\\$\frac{a}{b+c}+\frac{b}{a+c}+\frac{c}{a+b}\leq\frac{a^{\frac{3}{2}}+b^{\frac{3}{2}}+c^{\frac{3}{2}}}{2\sqrt{abc}}=\frac{1}{2\sqrt{abc}}$\\当且仅当$a=b=c$时取等号．
\end{solution}
\end{question}

